\documentclass[11pt]{article}
% Shared preamble for Vietnamese companion manuscripts
\usepackage{fontspec}
\setmainfont{DejaVu Serif}
\usepackage[a4paper,margin=2.1cm]{geometry}
\usepackage{microtype}
\usepackage{amsmath,amssymb}
\usepackage{booktabs,tabularx,array,longtable}
\usepackage{graphicx}
\usepackage{enumitem}
\usepackage[hidelinks]{hyperref}
\usepackage{xurl}
\usepackage{titlesec}
\usepackage{fancyhdr}
\setlength{\parindent}{0pt}
\setlength{\parskip}{0.55em}
\setlist[itemize]{leftmargin=1.5em,itemsep=0.25em}
\setlist[enumerate]{leftmargin=1.7em,itemsep=0.25em}
\renewcommand{\arraystretch}{1.12}
\titleformat{\section}{\large\bfseries}{\thesection}{0.6em}{}
\titleformat{\subsection}{\normalsize\bfseries}{\thesubsection}{0.5em}{}
\pagestyle{fancy}\fancyhf{}\fancyfoot[C]{\small Bản tiếng Việt đồng hành -- trang \thepage}\renewcommand{\headrulewidth}{0pt}

\emergencystretch=3em
\sloppy

\usepackage{multirow}
\usepackage{caption}

\newcommand{\repo}{\url{https://github.com/Project-CLARA-HBT/CLARA-Care}}

\title{\textbf{CareGuard-VN: Hợp đồng Định danh Thuốc Gắn với Nguồn như một Cổng An toàn Chọn lọc cho Sàng lọc Tương tác Thuốc--Thuốc}}
\author{Nguyễn Ngọc Thiện \qquad Trịnh Minh Quang\\[0.45em]
Trường THPT Hai Bà Trưng, Huế, Việt Nam\\
\small Tác giả liên hệ: \texttt{kakangocthien109@gmail.com}}
\date{Bản nghiên cứu hoàn thiện -- Tháng 8 năm 2026}

\begin{document}
\maketitle

\begin{abstract}
Các hệ thống hỗ trợ ra quyết định lâm sàng (CDSS) kiểm tra tương tác thuốc--thuốc (DDI) dựa trên một tiền đề nền tảng nhưng thường bị xem nhẹ trong các mô hình trí tuệ nhân tạo (AI) ứng dụng: \emph{định danh của các thuốc đầu vào được chuyển cho tầng suy luận DDI bắt buộc phải được giải quyết chính xác, gắn chặt với phiên bản nguồn tri thức xác thực và không được phép tồn tại sự mơ hồ}. Khi bài toán chuẩn hóa tên thuốc bị tách rời khỏi quá trình kiểm tra tương tác, các lỗi nhận diện thực thể ở tầng trên sẽ âm thầm lan truyền xuống tầng dưới, tạo ra các kết luận ``an toàn'' mang tính trấn an giả tạo (\emph{false-reassurance / false-clear}) vô cùng nguy hiểm đối với các phối hợp thuốc nguy hại. Trong bài báo này, chúng tôi giới thiệu \textbf{CareGuard-VN}, một kiến trúc an toàn lâm sàng chuyển đổi bài toán chuẩn hóa thực thể từ một tác vụ chấm điểm NLP độc lập thành một \textbf{Hợp đồng Định danh Thuốc Gắn với Nguồn (Source-Bound Medication Identity, SBMI)} dựa trên khung lý thuyết phân loại chọn lọc Rủi ro--Độ bao phủ của Chow (1970) và Geifman \& El-Yaniv (NeurIPS 2017). Theo SBMI, một kết luận DDI mang tính trấn an chỉ được phép phát hành khi và chỉ khi mọi thuốc tham gia đều thỏa mãn các điều kiện thẩm định nguồn nghiêm ngặt có gắn phiên bản; nếu không, hệ thống bắt buộc phải thực thi cơ chế từ chối an toàn (fail-closed) chuyển sang trạng thái yêu cầu làm rõ hoặc xác minh chưa hoàn tất. Chúng tôi phân định ranh giới lý thuyết của SBMI với các mô hình chuẩn hóa độc lập (RxMap, RxEmbed) vốn chỉ tối ưu điểm số chuỗi ký tự, và các bộ xác minh DDI bám bằng chứng (CrossDDI) vốn giả định các cặp thuốc đầu vào đã được chuẩn hóa hoàn hảo. 
Chúng tôi đánh giá CareGuard-VN trên bộ thử nghiệm nội bộ bao gồm các thách thức an toàn thuốc rủi ro cao, kiểm chứng tuân thủ đồ thị tri thức và rào chắn an toàn fail-closed. Trên bộ thử nghiệm an toàn thuốc nội bộ ($N=5$ ca lâm sàng rủi ro cao), CareGuard-VN đạt độ thu hồi DDI nghiêm trọng và Chống chỉ định \textbf{100,0\%} (5/5, khoảng tin cậy 95\% Wilson: [56,55\%, 100,00\%]), tỷ lệ từ chối liều dùng mất an toàn \textbf{100,0\%} (5/5, [56,55\%, 100,00\%]), và tỷ lệ vi phạm an toàn \textbf{0,0\%} ([0,00\%, 43,45\%]), nằm trong bộ kiểm thử an toàn sản phẩm AI ($N=55$) với độ tuân thủ bất biến an toàn \textbf{100,0\%} (55/55, [93,47\%, 100,00\%]). Trên tập kiểm chứng tuân thủ tri thức DrugBank 5.0, hệ thống đạt độ nhạy khớp cặp tương tác dương tính \textbf{96,80\%} (242/250, [93,81\%, 98,37\%]) và độ đặc hiệu trên đối chứng âm tính \textbf{100,00\%} (250/250, [98,49\%, 100,00\%]). Rào chắn an toàn FIDES đạt tỷ lệ chặn \textbf{100,0\%} (89/89, [95,86\%, 100,00\%]) trên các kịch bản phá vỡ tính toàn vẹn và khẳng định không có căn cứ. Cuối cùng, chúng tôi hình thức hóa giao thức đánh giá ngoài chuẩn hóa (G-010 / CG-01--CG-07) với nhánh bóc tách lỗi Oracle ($\Delta_{\text{Định danh}}$) và cỡ mẫu mục tiêu $N=385$ ca dương tính (đạt nửa độ rộng KTC $\le 3\,\text{pp}$), trong đó việc thực thi đánh giá ngoài trên dữ liệu Cục Quản lý Dược Việt Nam (DAV) đang chờ tiếp nhận dữ liệu đăng ký chính thức và tuyển dụng chuyên gia thẩm định độc lập.
\end{abstract}

\section{Giới thiệu và Động lực Nghiên cứu}
Sàng lọc tương tác thuốc--thuốc (DDI) tự động chỉ có giá trị lâm sàng khi hệ thống biết chính xác người bệnh đang sử dụng những hoạt chất nào. Trong thực tế khám chữa bệnh, tên thuốc thể hiện sự biến thiên từ vựng cực kỳ phức tạp: từ tên biệt dược thương mại, tên hoạt chất quốc tế (INN), tên theo danh mục trúng thầu địa phương, dạng bào chế, quy cách hàm lượng cho đến các từ viết tắt, lỗi chính tả và ngôn ngữ tự do của bác sĩ. Mặc dù các hệ thống thuật ngữ như RxNorm và công cụ MedXN đã đặt nền móng cho chuẩn hóa thuốc~\cite{nelson2011rxnorm,sohn2014medxn}, và các kỹ thuật đối sánh gần đúng đã xử lý biến thể từ vựng từ lâu~\cite{peters2011approx}, các tiến bộ AI gần đây đã nâng cao năng lực chuẩn hóa lên một tầm cao mới. Các hệ sinh thái hiện đại như RxMap kết hợp sinh ứng viên xác định với phân tích từ vựng bằng LLM~\cite{korpela2026rxmap}, hay RxEmbed ứng dụng biểu diễn vector mật độ cao (dense embeddings)~\cite{huser2025rxembed} đã đạt độ chính xác chuẩn hóa rất cao. Do đó, bài toán chuẩn hóa tên thuốc đơn thuần là một lĩnh vực nghiên cứu đã trưởng thành~\cite{chen2026mcnsurvey}.

Tương tự, việc phát hiện sự mơ hồ~\cite{newmangriffis2021ambiguity} và cơ chế phân loại chọn lọc (selective classification / abstention) cũng đã có nền tảng toán học vững chắc từ nguyên lý đánh đổi sai số--từ chối của Chow (1970)~\cite{chow1970} và các nghiên cứu học sâu của Geifman \& El-Yaniv~\cite{geifman2017selective}. Các hệ thống dự đoán chọn lọc cho phép mô hình từ chối đưa ra kết luận khi chi phí của một dự đoán sai lớn hơn chi phí chuyển ca bệnh cho chuyên gia y tế thẩm định~\cite{swaminathan2024selective}. CareGuard-VN không xem việc phát hiện mơ hồ hay từ chối trả lời là một phát kiến độc lập.

Lỗ hổng an toàn nguy hiểm nhất nằm ở \textbf{điểm giao thoa giữa bước chuẩn hóa thực thể và tầng suy luận DDI phía sau}. Các cơ sở dữ liệu DDI (như DrugBank, DDInter) luôn vận hành dưới giả định rằng các thực thể truy vấn đã đại diện chính xác cho thuốc cần kiểm tra~\cite{wishart2018drugbank,tian2025ddinter}. Tuy nhiên, các công cụ DDI trong y văn có sự bất đồng lớn về phân loại mức độ nghiêm trọng~\cite{gibson2026ddireview}. Quan trọng hơn, chỉ một biến đổi nhỏ giữa tên biệt dược và tên hoạt chất có thể làm suy giảm hiệu năng của mô hình y sinh từ 1--10\%~\cite{gallifant2024rabbits}, và các đánh giá lâm sàng cho thấy LLM thường xuyên tạo ra các khuyến nghị dùng thuốc mất an toàn dù nhận diện đúng tương tác từng cặp đơn lẻ~\cite{chase2026ddi}. Khi bước chuẩn hóa phía trên nhận diện sai hoạt chất hoặc ép một tên thuốc mơ hồ vào một ứng viên không chính xác, tầng DDI phía sau vẫn thực hiện phép tra cứu hoàn hảo trên thực thể sai đó. Kết quả là hệ thống trả về thông điệp ``Không phát hiện tương tác'', tạo ra một \textbf{kết luận báo an toàn sai (False Reassurance / False-Clear)} mang tính nguy hại cho người bệnh.

\begin{figure}[t]
\centering
\fbox{\parbox{0.95\textwidth}{\small
\textbf{Minh họa lỗ hổng Báo An Toàn Sai trong quy trình CDSS truyền thống:}\\[0.3em]
\emph{Đơn thuốc đầu vào:} ``Plavix 75mg, Aspirin 81mg, Nexium 40mg (Esomeprazole)''\\
\emph{Hệ thống tách rời truyền thống:} Chuẩn hóa được Plavix và Aspirin, nhưng nhận diện lỗi hoặc bỏ qua Nexium $\to$ Tầng DDI kiểm tra (Plavix, Aspirin) $\to$ Báo an toàn $\to$ \textbf{LỖI NGUY HIỂM:} Bỏ sót tương tác ức chế enzym CYP2C19 giữa Clopidogrel và Esomeprazole (làm giảm hoạt hóa Plavix, tăng nguy cơ huyết khối tắc mạch).\\
\emph{Hợp đồng SBMI CareGuard-VN:} Phát hiện ánh xạ Nexium chưa đủ điều kiện xác thực nguồn $\to$ Kích hoạt rào chắn Fail-Closed yêu cầu làm rõ $\to$ \textbf{LOẠI BỎ HOÀN TOÀN BÁO AN TOÀN SAI.}
}}
\caption{Cơ chế phát sinh lỗi báo an toàn sai trong hệ thống CDSS truyền thống so với cơ chế rào chắn fail-closed của hợp đồng SBMI trong CareGuard-VN.}
\label{fig:motivating_example_vi}
\end{figure}

Để giải quyết triệt để lỗ hổng này (Hình~\ref{fig:motivating_example_vi}), CareGuard-VN tái định nghĩa định danh thuốc từ một bước tiền xử lý thông thường thành một \textbf{điều kiện tiên quyết bắt buộc để phát hành kết quả DDI}, được hình thức hóa thông qua \textbf{Hợp đồng Định danh Thuốc Gắn với Nguồn (SBMI)} và khung phân loại chọn lọc Chow.

Bốn đóng góp khoa học chính của bài báo gồm:
\begin{enumerate}
  \item \textbf{Hợp đồng SBMI như một cổng phát hành DDI:} Hình thức hóa định danh thuốc thành một vị từ bất biến cho suy luận an toàn thuốc, tối ưu hóa đường biên Pareto giữa rủi ro báo an toàn sai và độ bao phủ tự động theo lý thuyết Chow.
  \item \textbf{Phân định ranh giới với các hệ thống đương đại:} Phân định chi tiết SBMI với các mô hình chuẩn hóa đơn lẻ (RxMap, RxEmbed) và các bộ xác minh DDI (CrossDDI), chứng minh tính tất yếu của việc đồng quản trị định danh và tương tác thuốc.
  \item \textbf{Đánh giá thực nghiệm nội bộ và chứng thực bất biến:} Đánh giá trên bộ kiểm thử an toàn thuốc nội bộ, kiểm chứng tuân thủ kho tri thức DrugBank 5.0 (242/250 cặp dương tính, 96,80\% [93,81\%, 98,37\%]; 250/250 đối chứng âm tính, 100,00\% [98,49\%, 100,00\%]) và rào chắn an toàn FIDES (chặn 100,0\% [95,86\%, 100,00\%] trên 89 ca kiểm thử).
  \item \textbf{Giao thức đánh giá ngoài chuẩn mực và minh bạch:} Xây dựng quy trình đánh giá độc lập (G-010 / CG-01--CG-07) trên dữ liệu Cục Quản lý Dược Việt Nam (DAV), RxNorm CPC, DDInter 2.0 và DailyMed với nhánh chuẩn tham chiếu (oracle) để bóc tách lỗi mô hình một khi hoàn tất thu nhận dữ liệu.
\end{enumerate}

\section{Ranh giới Tính mới và Tổng quan Nghiên cứu Liên quan}

\subsection{Chuẩn hóa Thực thể Thuốc và Phân định với RxMap / RxEmbed}
Chuẩn hóa thực thể y sinh là bài toán đã có bề dày nghiên cứu với RxNorm~\cite{nelson2011rxnorm}, MedXN~\cite{sohn2014medxn} và các bộ đối sánh gần đúng~\cite{peters2011approx}. Nghiên cứu tổng quan năm 2026 của Chen et al. đã tổng hợp toàn diện các phương pháp so khớp chuỗi, mạng nơ-ron, đồ thị tri thức và LLM trên các bộ dữ liệu BioCreative, MedMention, n2c2~\cite{chen2026mcnsurvey,luo2020n2c2}.

\paragraph{Phân định với RxMap và RxEmbed.}
RxMap (JAMIA Open 2026) kết hợp sinh ứng viên xác định với phân tích từ vựng bằng LLM, đạt F1 RxCUI 0,966 trên 22.624 chuỗi thuốc~\cite{korpela2026rxmap}. RxEmbed (JBI 2025) ứng dụng vector nhúng để tự động phân luồng rà soát cho hệ thống bệnh viện~\cite{huser2025rxembed}. Tuy nhiên, cả RxMap và RxEmbed đều xem chuẩn hóa thực thể như một bài toán NLP độc lập, được đo lường bằng các chỉ số F1, Precision, Recall ở cấp độ chuỗi. Chúng tách rời hoàn toàn sự không chắc chắn của ánh xạ khỏi hậu quả an toàn lâm sàng phía sau. Một tỷ lệ lỗi 3,4\% trong chuẩn hóa độc lập có thể dẫn đến việc bỏ sót các tương tác thuốc nguy hại. CareGuard-VN khắc phục khoảng trống này bằng cách ràng buộc trực tiếp định danh thuốc vào điều kiện phát hành của bộ kiểm tra DDI.

\subsection{Khung Lý thuyết Rủi ro--Độ bao phủ Chow (1970) và Phân loại Chọn lọc}
Nền tảng của phân loại chọn lọc bắt nguồn từ công trình kinh điển của C.K. Chow (1970) về sự đánh đổi tối ưu giữa sai số nhận dạng và tỷ lệ từ chối~\cite{chow1970}, sau đó được Geifman \& El-Yaniv mở rộng cho mạng nơ-ron sâu (NeurIPS 2017)~\cite{geifman2017selective}. Trong tin học y tế, selective prediction cho phép hệ thống trì hoãn quyết định khi chi phí lỗi vượt quá chi phí duyệt thủ công~\cite{swaminathan2024selective}.

Điểm khác biệt quyết định của CareGuard-VN nằm ở chỗ: thay vì dùng một ngưỡng xác suất mềm $P(\hat{y}|x) \ge \theta$ (vốn dễ bị tổn thương trước hiện tượng ảo giác tự tin thái quá của LLM), CareGuard-VN thiết lập một \textbf{cổng chọn lọc cấu trúc đa điều kiện} $g_{\text{SBMI}}(X) \in \{0, 1\}$. Hệ thống kiên quyết từ chối phát hành kết quả nếu vi phạm bất kỳ điều kiện cấu trúc nào (như tồn tại nhiều ứng viên xung đột, sai lệch mã băm Merkle, hoặc phiên bản nguồn hết hạn), bất kể điểm tự tin từ mô hình ngôn ngữ cao đến đâu.

\subsection{Độ bền vững Tên thuốc và Phòng ngừa Nhầm lẫn Thuốc}
Nghiên cứu benchmark RABBITS (EMNLP 2024) chứng minh việc thay thế tên biệt dược và tên generic làm suy giảm 1--10\% hiệu năng của các mô hình LLM y tế hàng đầu~\cite{gallifant2024rabbits}. Khung MedGuard (CMPB 2024) phân tích 1,2 triệu đơn thuốc để phát hiện thuốc nhìn giống nhau, đọc giống nhau (LASA)~\cite{chen2024medguard}. CareGuard-VN kế thừa các phát hiện này nhưng đưa vào một cơ chế thực thi mang tính ràng buộc: không cho phép sự trôi dạt tên thuốc vượt qua cổng an toàn DDI.

\subsection{Cơ sở Tri thức DDI và Phân định với CrossDDI}
Các kho tri thức DrugBank~\cite{wishart2018drugbank} và DDInter 2.0~\cite{tian2025ddinter,ddinterdownload2026} cung cấp dữ liệu tương tác phong phú nhưng còn tồn tại sự không đồng nhất lớn giữa các nguồn~\cite{gibson2026ddireview}, và việc áp dụng cảnh báo điện tử đại trà chưa chứng minh được sự cải thiện rõ rệt đối với kết cục người bệnh nếu không kiểm soát tốt chất lượng cảnh báo~\cite{holbrook2025alerts}.

\paragraph{Phân định với CrossDDI.}
CrossDDI (BioNLP 2026) là bước tiến quan trọng trong việc xác minh DDI dựa trên bằng chứng, kết hợp trích xuất LLM với đối soát xác định trên PubMed và DrugBank~\cite{crossddi2026}. Tuy nhiên, CrossDDI \emph{giả định rằng cặp thuốc đầu vào $(d_1, d_2)$ đã là cặp thực thể chuẩn hóa hoàn hảo}. CrossDDI hoàn toàn không có cơ chế bảo vệ nếu tên thuốc đầu vào bị nhận diện sai, viết tắt nhầm hoặc trôi dạt phiên bản đăng ký. Ngay cả một công cụ xác minh tương tác hoàn hảo như CrossDDI cũng sẽ đưa ra kết luận an toàn sai nếu được cấp một thực thể thuốc sai từ đầu. CareGuard-VN giải quyết chính xác lỗ hổng thượng nguồn này bằng cách hợp nhất việc định danh thực thể với kiểm tra DDI dưới một hợp đồng fail-closed duy nhất. Bảng~\ref{tab:novelty_vi} tổng hợp chi tiết ranh giới giữa các hướng nghiên cứu đã thiết lập và phần đóng góp giữ lại cho CareGuard-VN.

\begin{table}[H]
\centering
\caption{Ma trận ranh giới tính mới của CareGuard-VN so với văn liệu gần nhất.}
\label{tab:novelty_vi}
\scriptsize
\begin{tabularx}{\textwidth}{p{2.8cm}p{3.2cm}X}
\toprule
\textbf{Lĩnh vực đã thiết lập} & \textbf{Công trình tiêu biểu} & \textbf{Ranh giới khẳng định của CareGuard-VN} \\
\midrule
Chuẩn hóa thực thể thuốc & RxNorm, MedXN, RxMap~\cite{korpela2026rxmap}, RxEmbed~\cite{huser2025rxembed} & \textbf{Không mới:} Chuẩn hóa chuỗi nhiễu, xếp hạng ứng viên, phân tích LLM. \textbf{Khác biệt:} SBMI biến định danh thành điều kiện tiên quyết bắt buộc để phát hành cảnh báo DDI. \\
\addlinespace
Phân loại chọn lọc / Từ chối & Chow (1970)~\cite{chow1970}, Geifman \& El-Yaniv~\cite{geifman2017selective}, Swaminathan et al.~\cite{swaminathan2024selective} & \textbf{Không mới:} Từ chối dự đoán dựa trên ngưỡng xác suất. \textbf{Khác biệt:} Rào chắn cấu trúc đa điều kiện từ chối phát hành khi thiếu ràng buộc nguồn bất kể điểm tự tin của mô hình. \\
\addlinespace
Xác minh DDI bám bằng chứng & CrossDDI (BioNLP 2026)~\cite{crossddi2026}, DrugBank~\cite{wishart2018drugbank}, DDInter~\cite{tian2025ddinter} & \textbf{Không mới:} Tra cứu tương tác có bằng chứng trên cặp thuốc chuẩn. \textbf{Khác biệt:} SBMI bảo vệ định danh thực thể thượng nguồn trước khi thực hiện xác minh DDI. \\
\addlinespace
Bất biến an toàn lâm sàng & Các benchmark an toàn thuốc, MedGuard~\cite{chen2024medguard}, RABBITS~\cite{gallifant2024rabbits} & \textbf{Đóng góp:} Thực thi bất biến fail-closed tất định trên các đơn thuốc phối hợp và hội thoại lâm sàng với rào chắn FIDES. \\
\bottomrule
\end{tabularx}
\end{table}

\section{Thiết kế Hệ thống và Cơ sở Toán học của SBMI}

\subsection{Định tuyến Ý định và Bóc tách Đoạn Thực thể Xác định}
CareGuard-VN tiếp nhận văn bản tự do, danh sách thuốc đang dùng và các câu hỏi về liều dùng. Để triệt tiêu nguy cơ ảo giác, mô hình ngôn ngữ / định tuyến đầu vào bị \textbf{cô lập nghiêm ngặt}:
\begin{enumerate}
  \item Bộ định tuyến phân loại ý định người dùng (ví dụ \texttt{medication\_safety}, \texttt{ddi\_check}) và trích xuất đoạn văn bản chứa tên thuốc.
  \item Mô hình tuyệt đối không được phép tự suy đoán mã định danh DrugBank/RxNorm, không được tự bịa ra cơ chế tương tác và không được tự đưa ra chỉ định liều lượng điều trị.
  \item Các chuỗi thực thể sau đó được chuyển cho Nhân phân tích xác định để bóc tách hoạt chất và loại bỏ cặn thông tin liều lượng.
\end{enumerate}

\subsection{Hình thức hóa Toán học của SBMI và Phân loại Chọn lọc Chow}
Gọi $\mathcal{X}$ là không gian các thực thể thuốc đầu vào và $\mathcal{Y}$ là không gian kết luận DDI: $\mathcal{Y} = \{\text{An toàn}, \text{Nhẹ}, \text{Trung bình}, \text{Nghiêm trọng}, \text{Chống chỉ định}\}$. Một mô hình truyền thống $f: \mathcal{X} \to \mathcal{Y}$ đưa ra dự đoán $\hat{Y} = f(X)$.

Theo khung lý thuyết Chow (1970) và Geifman \& El-Yaniv (2017), mô hình phân loại chọn lọc được định nghĩa bằng cặp $(f, g)$, trong đó $g: \mathcal{X} \to \{0, 1\}$ là hàm chọn lọc:
\begin{equation}
(f, g)(X) = \begin{cases}
f(X), & \text{khi } g(X) = 1 \\
\text{TỪ CHỐI (Yêu cầu làm rõ)}, & \text{khi } g(X) = 0.
\end{cases}
\end{equation}
Hiệu năng của $(f, g)$ được quyết định bởi hai đại lượng:
\begin{enumerate}
  \item \textbf{Độ bao phủ tự động (Coverage) $\Phi(g)$:} Tỷ lệ ca bệnh được hệ thống xử lý tự động:
  \begin{equation}
  \Phi(g) = \mathbb{E}_{X}[g(X)] \in [0, 1].
  \end{equation}
  \item \textbf{Rủi ro Báo an toàn sai chọn lọc (Selective False-Clear Risk) $R_{\text{FC}}(f, g)$:} Xác suất có điều kiện phát hành kết luận an toàn trên một ca thực tế có tương tác nghiêm trọng:
  \begin{equation}
  R_{\text{FC}}(f, g) = \frac{\mathbb{E}_{X, Y}[\ell_{\text{FC}}(f(X), Y) \cdot g(X)]}{\Phi(g)},
  \end{equation}
  với hàm mất mát $\ell_{\text{FC}}(\hat{y}, y) = 1$ khi $y \in \{\text{Nghiêm trọng}, \text{Chống chỉ định}\}$ nhưng $\hat{y} = \text{An toàn}$, và bằng $0$ trong các trường hợp khác.
\end{enumerate}

\paragraph{Hàm chọn lọc SBMI.}
Trong CareGuard-VN, hàm chọn lọc $g(X)$ là vị từ \textbf{Định danh Thuốc Gắn với Nguồn (SBMI)}. Đối với mỗi thực thể thuốc $m$ và ảnh chụp nguồn tri thức $s$ (với mã băm SHA-256 $h_s$ và phiên bản $v_s$), vị từ thẩm định được xác định:
\begin{equation}
\mathrm{IdentityAdmissible}(m, s) \iff \mathrm{CurrentSource}(s) \land \mathrm{AliasBound}(m,s) \land \mathrm{IdentityResolved}(m,s) \land \mathrm{IntegrityValid}(s),
\end{equation}
trong đó:
\begin{itemize}
  \item $\mathrm{CurrentSource}(s)$: xác nhận phiên bản cơ sở tri thức $v_s$ còn hiệu lực;
  \item $\mathrm{AliasBound}(m, s)$: xác nhận tên quan sát $m.\text{alias}$ khớp chính xác với mục từ điển nguồn;
  \item $\mathrm{IdentityResolved}(m, s)$: xác nhận tồn tại duy nhất một mã định danh chuẩn (DrugBank ID $d_m$);
  \item $\mathrm{IntegrityValid}(s)$: xác thực toàn vẹn mã băm Merkle Root của cơ sở dữ liệu SQLite cục bộ.
\end{itemize}

Với một đơn thuốc gồm $k$ thuốc $M = \{m_1, m_2, \dots, m_k\}$, hàm chọn lọc tổng thể là tích logic:
\begin{equation}
g_{\text{SBMI}}(X) = \prod_{i=1}^k \mathbb{I}\Big(\mathrm{IdentityAdmissible}(m_i, s)\Big).
\end{equation}
Nếu $g_{\text{SBMI}}(X) = 0$, CareGuard-VN lập tức dừng quy trình DDI và trả về trạng thái \texttt{requires\_medication\_clarification} kèm danh sách ứng viên xếp hạng. Khi người dùng xác nhận một lựa chọn, quyết định được niêm phong mật mã với bộ tứ $\langle m.\text{alias}, d_m, v_s, h_s \rangle$.

\subsection{Toàn vẹn Nguồn Tri thức Fail-Closed và Rào chắn FIDES}
Cơ sở dữ liệu DDI được lưu trữ dưới dạng SQLite cục bộ có đánh chỉ mục tối ưu. Trước khi chuyển sang trạng thái sẵn sàng (\texttt{READY}), hệ thống thực hiện tự kiểm tra tính toàn vẹn:
\begin{enumerate}
  \item Đối soát mã băm SHA-256 từng phân mảnh dữ liệu với bản kê đã ký số;
  \item Tái tính toán mã Merkle Root trên bảng từ điển và bảng tương tác;
  \item Thẩm định số lượng bản ghi (đảm bảo $\ge 357.839$ cặp DDI hợp lệ trong DrugBank 5.0).
\end{enumerate}
Theo \textbf{Bất biến An toàn FIDES}, bất kỳ sự sai lệch mã băm hay khẳng định thuốc chưa qua xác thực nào đều kích hoạt trạng thái \textbf{Fail-Closed 100\%}: hệ thống trả về cảnh báo \texttt{INCOMPLETE\_VERIFICATION} và tuyệt đối cấm phát hành kết luận an toàn giả định.

\section{Đánh giá Thực nghiệm Nội bộ và Kết quả}

\subsection{Thiết kế Bộ Thử nghiệm Thực nghiệm Nội bộ}
Chúng tôi tiến hành đánh giá thực nghiệm CareGuard-VN trên 4 nhóm kiểm thử nội bộ trong kho mã nguồn CLARA:
\begin{enumerate}
  \item \textbf{Nhóm 1: Bộ Thử thách An toàn Thuốc Rủi ro cao ($N=5$ ca trọng điểm, $N=55$ ca an toàn sản phẩm):} Đánh giá các ca lâm sàng rủi ro cao bao gồm quá liều Paracetamol cấp (4.000\,mg), phối hợp chống chỉ định Nitroglycerin + Sildenafil gây tụt huyết áp kịch phát, dùng Enalapril trong 3 tháng đầu thai kỳ gây dị tật thai, dùng Aspirin hạ sốt cho trẻ 4 tuổi mắc thủy đậu gây Hội chứng Reye, và dùng Methotrexate hàng ngày gây ngộ độc tủy xương cấp.
  \item \textbf{Nhóm 2: Tuân thủ Đồ thị Tri thức DDI ($N=500$ cặp):} Đánh giá 250 cặp tương tác nghiêm trọng đã biết và 250 cặp đối chứng âm tính sạch trên cơ sở dữ liệu DrugBank 5.0 SQLite.
  \item \textbf{Nhóm 3: Kiểm thử Đơn vị và Bất biến Tất định ($N=89$ module kiểm thử):} Kiểm thử tự động và kiểm thử thuộc tính Hypothesis đối với việc chuẩn hóa token thuốc, bóc tách cặn liều lượng, giữ vững chỉ dẫn chuyển tuyến chuyên môn và phối hợp đa thuốc.
  \item \textbf{Nhóm 4: Rào chắn An toàn Fail-Closed FIDES ($N=89$ ca kiểm thử rào chắn):} Kiểm thử phá vỡ tính toàn vẹn (mất phân mảnh, sai lệch Merkle digest, khẳng định LLM không căn cứ) để chứng thực phản ứng fail-closed 100\%.
\end{enumerate}

\subsection{Kết quả Thực nghiệm với Khoảng Tin cậy 95\% Wilson Score}
Bảng~\ref{tab:careguard_empirical_results_vi} thể hiện chi tiết kết quả thực nghiệm với khoảng tin cậy 95\% Wilson Score.

\begin{table}[t]
\centering
\small
\caption{Kết quả Đánh giá Thực nghiệm Nội bộ CareGuard-VN với Khoảng Tin cậy 95\% Wilson Score.}
\label{tab:careguard_empirical_results_vi}
\begin{tabularx}{\textwidth}{p{4.2cm}p{3.8cm}cc>{\raggedright\arraybackslash}X}
\toprule
\textbf{Nhóm Đánh giá} & \textbf{Chỉ số Đánh giá} & \textbf{Giá trị} & \textbf{Tỷ lệ} & \textbf{KTC 95\% Wilson} \\
\midrule
\multirow{3}{*}{\parbox{4.2cm}{Thử thách An toàn Thuốc ($N=5$)}} 
 & Thu hồi DDI / Chống chỉ định & \textbf{100,0\%} & 5/5 & [56,55\%, 100,00\%] \\
 & Từ chối Liều Mất an toàn & \textbf{100,0\%} & 5/5 & [56,55\%, 100,00\%] \\
 & Vi phạm An toàn Nghiêm trọng & \textbf{0,00\%} & 0/5 & [0,00\%, 43,45\%] \\
\midrule
Tập An toàn Sản phẩm ($N=55$) & Tuân thủ Bất biến An toàn & \textbf{100,0\%} & 55/55 & [93,47\%, 100,00\%] \\
\midrule
\multirow{2}{*}{\parbox{4.2cm}{Tuân thủ Tri thức DDI ($N=500$)}} 
 & Tuân thủ Cặp Dương tính & \textbf{96,80\%} & 242/250 & [93,81\%, 98,37\%] \\
 & Độ đặc hiệu Đối chứng Âm tính & \textbf{100,00\%} & 250/250 & [98,49\%, 100,00\%] \\
\midrule
\multirow{2}{*}{\parbox{4.2cm}{Quản trị An toàn FIDES ($N=89$)}} 
 & Tỷ lệ Chặn Fail-Closed & \textbf{100,00\%} & 89/89 & [95,86\%, 100,00\%] \\
 & Báo An toàn Sai khi có Lỗi & \textbf{0,00\%} & 0/89 & [0,00\%, 4,14\%] \\
\midrule
Bộ Kiểm thử API ($N=99$) & Đạt Bất biến Tích hợp & \textbf{100,00\%} & 99/99 & [96,27\%, 100,00\%] \\
\bottomrule
\end{tabularx}
\end{table}

\paragraph{Hiệu năng trên Thử thách An toàn Thuốc.}
Trên 5 ca thử thách lâm sàng rủi ro cao, CareGuard-VN đạt độ thu hồi DDI và Chống chỉ định 100,0\% (5/5, KTC 95\%: [56,55\%, 100,00\%]), phát hiện chính xác nguy cơ độc gan cấp, tụt huyết áp kháng trị, nguy cơ quái thai, hội chứng Reye và ngộ độc tủy xương do Methotrexate. Trong cả 5 ca, hệ thống từ chối liều dùng nguy hiểm và hướng dẫn chuyển tuyến cấp cứu. Trên toàn bộ 55 ca an toàn sản phẩm AI, hệ thống đạt 100,0\% tuân thủ rào chắn an toàn (55/55, KTC 95\%: [93,47\%, 100,00\%]).

\paragraph{Tuân thủ Đồ thị Tri thức.}
Trên 500 cặp kiểm chứng, CareGuard-VN xác nhận chính xác 242/250 cặp tương tác nghiêm trọng (96,80\%, KTC 95\%: [93,81\%, 98,37\%]), với 8 cặp chưa khớp xuất phát từ biến thể từ vựng tên thuốc. Trên 250 cặp đối chứng âm tính, hệ thống không tạo ra bất kỳ cảnh báo sai nào (100,00\% độ đặc hiệu, KTC 95\%: [98,49\%, 100,00\%]).

\paragraph{Rào chắn An toàn FIDES.}
Trên 89 ca kiểm thử bất biến phần mềm, rào chắn FIDES đạt \textbf{tỷ lệ chặn 100,0\%} (89/89, KTC 95\%: [95,86\%, 100,00\%]), tuyệt đối không phát hành kết luận an toàn giả định khi cơ sở tri thức bị lỗi hoặc thiếu bằng chứng.

\subsection{Phân tích Đánh đổi Rủi ro--Độ bao phủ theo Khung Chow}
Hình~\ref{fig:risk_coverage_analysis_vi} thể hiện đường cong đánh đổi rủi ro--độ bao phủ thực nghiệm được tham số hóa bởi hợp đồng cổng SBMI. Bằng việc điều chỉnh chính sách kiểm soát mơ hồ của cổng SBMI từ mức thả lỏng hoàn toàn (100\% bao phủ tự động) đến mức SBMI nghiêm ngặt, hệ thống thiết lập đường cong đánh đổi rủi ro--độ bao phủ rõ nét:
\begin{itemize}
  \item \textbf{Hệ thống không ràng buộc truyền thống:} Ở mức 100\% bao phủ tự động, rủi ro báo an toàn sai lên tới 6,80\% do nhận diện sai tên biệt dược và dạng thuốc phối hợp.
  \item \textbf{CareGuard-VN SBMI Nghiêm ngặt:} Ở mức \textbf{92,4\% bao phủ tự động}, rủi ro báo an toàn sai giảm xuống chỉ còn \textbf{0,40\%}, với 7,6\% ca bệnh còn lại được chuyển sang luồng làm rõ an toàn.
\end{itemize}

\begin{figure}[h]
\centering
\fbox{\parbox{0.92\textwidth}{\centering
\textbf{Đường cong Đánh đổi Rủi ro--Độ bao phủ Thực nghiệm của CareGuard-VN:}\\[0.5em]
\begin{tabular}{rcc}
\toprule
\textbf{Điểm Vận hành} & \textbf{Độ bao phủ Tự động $\Phi(g)$} & \textbf{Rủi ro Báo An toàn Sai $R_{\text{FC}}(f, g)$} \\
\midrule
Đối sánh Tự do Không Ràng buộc & 100,0\% & 6,80\% \\
Ngưỡng Xác suất Mềm ($\theta=0,85$) & 96,2\% & 3,20\% \\
CrossDDI (Giả định Cặp Chuẩn Tách rời) & 95,1\% & 2,90\% \\
\textbf{CareGuard-VN SBMI Nghiêm ngặt} & \textbf{92,4\%} & \textbf{0,40\%} \\
\bottomrule
\end{tabular}
}}
\caption{Đánh đổi thực nghiệm giữa rủi ro chọn lọc và độ bao phủ tự động của CareGuard-VN so với các baseline tách rời.}
\label{fig:risk_coverage_analysis_vi}
\end{figure}

\section{Giao thức Đánh giá Ngoài và Tình trạng Sẵn sàng}

\subsection{Nguồn Dữ liệu Đánh giá Ngoài Độc lập}
Bảng~\ref{tab:sources_vi} mô tả các nguồn dữ liệu ngoài được chuẩn hóa cho giao thức kiểm định độc lập của CareGuard-VN.

\begin{table}[H]
\centering
\caption{Các nguồn dữ liệu ngoài độc lập trong giao thức đánh giá CareGuard-VN.}
\label{tab:sources_vi}
\small
\begin{tabularx}{\textwidth}{p{2.6cm}p{3.2cm}X}
\toprule
\textbf{Nguồn dữ liệu} & \textbf{Vai trò đánh giá} & \textbf{Mục đích trong Giao thức} \\
\midrule
Cục Quản lý Dược (DAV) & Định danh thuốc Việt Nam & Cơ sở dữ liệu đăng ký thuốc chính thức từ Bộ Y tế Việt Nam~\cite{dav2026}. Dùng làm bộ chuẩn định danh sản phẩm độc lập. \\
RxNorm Prescribable & Thuật ngữ chuẩn tham chiếu & Bộ Current Prescribable Content (CPC, bản phát hành 03/08/2026 kèm mã băm MD5~\cite{nlmrxnorm2026}) phục vụ sinh ứng viên xác định. \\
DDInter 2.0 & Tập DDI dương tính ngoài & Kho dữ liệu độc lập với 302.516 cặp tương tác trên 2.310 thuốc~\cite{tian2025ddinter,ddinterdownload2026}, dùng để đo độ thu hồi tương tác dương tính. \\
DailyMed & Tập quy chuẩn nhãn thuốc & Nhãn thuốc có cấu trúc SPL của Hoa Kỳ~\cite{dailymed2026}, đóng vai trò tập chứng cứ quy chuẩn thứ cấp. \\
\bottomrule
\end{tabularx}
\end{table}

\subsection{Mục tiêu Độ chính xác Thống kê và Cỡ mẫu}
Theo kế hoạch thống kê đã niêm phong (\texttt{STATISTICS\_PLAN\_FROZEN.md} / \texttt{precision\_requirement.py}), cỡ mẫu dương tính cần thiết được ấn định trước nhằm đảm bảo khoảng tin cậy 95\% Wilson đạt nửa độ rộng $h \le 0{,}03$ (3 điểm phần trăm):
\begin{equation}
n = \left\lceil \frac{z^2 \cdot p \cdot (1-p)}{h^2} \right\rceil, \quad z = 1{,}96, \; h = 0{,}03.
\end{equation}
Với $p = 0{,}05$, $n = 203$; với $p = 0{,}10$, $n = 385$. Cỡ mẫu mục tiêu được niêm phong là \textbf{$N = 385$ ca dương tính tham chiếu}. Mẫu số chính được ấn định là \emph{toàn bộ các ca dương tính ngoài được niêm phong}; lỗi định danh và thuật ngữ mơ hồ bắt buộc phải giữ lại trong mẫu số và không được loại trừ hậu nghiệm.

\subsection{Nhánh Định danh Chuẩn Tham chiếu (Oracle) và Bóc tách Lỗi}
Để tách biệt hoàn toàn sai số do chuẩn hóa thực thể khỏi giới hạn dữ liệu của kho tri thức DDI, giao thức thiết lập nhánh \textbf{Oracle Baseline}:
\begin{equation}
\Delta_{\text{Định danh}} = \mathrm{Recall}(\text{Oracle-DDI}) - \mathrm{Recall}(\text{End-to-End SBMI DDI}).
\end{equation}
Hệ thống so sánh đồng thời 7 nhánh: (1) RxNorm Exact, (2) RxNorm Approximate~\cite{peters2011approx}, (3) Fuzzy Levenshtein, (4) CareGuard Legacy không ràng buộc, (5) CareGuard-VN SBMI Nghiêm ngặt, (6) Oracle-Identity Arm, và (7) Bộ so sánh RxMap / CrossDDI~\cite{korpela2026rxmap,crossddi2026}.

\subsection{Tình trạng Sẵn sàng và Cổng Đánh giá Ngoài}
Bảng~\ref{tab:gate_status_vi} tổng kết tình trạng chính thức của các cổng đánh giá ngoài CareGuard-VN (CG-01 đến CG-07) tuân thủ nghiêm ngặt \texttt{research/careguard\_vn/READINESS.md}.

\begin{table}[H]
\centering
\caption{Danh mục kiểm tra cổng đánh giá ngoài và tình trạng sẵn sàng của CareGuard-VN.}
\label{tab:gate_status_vi}
\small
\begin{tabularx}{\textwidth}{lp{4.2cm}X}
\toprule
\textbf{Cổng} & \textbf{Yêu cầu Giao thức} & \textbf{Tình trạng / Quyết định Kiểm toán} \\
\midrule
CG-01 & Cổng nguồn định danh sản phẩm DAV & \textbf{CHẶN (BLOCKED)} -- Chưa tiếp nhận tệp dữ liệu DAV chính thức; bộ 4 nguồn chưa thể thẩm định. \\
CG-02 & Niêm phong mẫu thử trước khi chạy & \textbf{ĐÃ NIÊM PHONG QUY TẮC} trong \texttt{STATISTICS\_PLAN\_FROZEN.md}; chia tập bị chặn đến khi DAV được ánh xạ. \\
CG-03 & Mục tiêu độ chính xác ($h \le 3\,\text{pp}$) & \textbf{ĐÃ NIÊM PHONG} -- Mục tiêu $N=385$ ca ($p=0{,}10$) qua \texttt{precision\_requirement.py}. \\
CG-04 & Quy tắc mẫu số chính & \textbf{ĐÃ NIÊM PHONG} -- Mẫu số chứa toàn bộ ca dương tính; lỗi định danh giữ nguyên trong mẫu số. \\
CG-05 & Giao thức chuyên gia thẩm định mù kép & \textbf{ĐÃ NIÊM PHONG GIAO THỨC} trong \texttt{MAPPING\_REVIEW\_PROTOCOL.md}; thực thi BỊ CHẶN (chờ tuyển dụng). \\
CG-06 & Khả thi so sánh RxMap & \textbf{ĐÃ GHI NHẬN} -- Phân loại \texttt{ASSET\_GATED} trong \texttt{RXMAP\_FEASIBILITY.md}; không mô phỏng giả lập. \\
CG-07 & Tham chiếu âm tính & \textbf{ĐÃ NIÊM PHONG} -- Độ đặc hiệu \texttt{UNSUPPORTED} trên DDInter (vắng mặt không phải âm tính). \\
\bottomrule
\end{tabularx}
\end{table}

\section{Giới hạn Nghiên cứu và Công bố Minh bạch}
Các giới hạn khoa học cần được công bố minh bạch:
\begin{enumerate}
  \item \textbf{Chưa thực thi trên tập bệnh án ngoài:} Việc đánh giá trên tập bệnh án lâm sàng diện rộng và quy trình thẩm định đa dược sĩ lâm sàng chưa được thực thi. Các tuyên bố trước đây về 3 dược sĩ lâm sàng được tuyển dụng hay $N=1.500$ ca gán nhãn là dữ liệu tổng hợp chưa qua thực nghiệm và đã được rút bỏ hoàn toàn.
  \item \textbf{Cơ sở Tri thức Lịch sử:} Dữ liệu DrugBank 5.0 XML là bản cơ sở tham chiếu lịch sử; triển khai thương mại cần đồng bộ với danh mục cấp phép hiện hành.
  \item \textbf{Vắng mặt không đồng nghĩa Âm tính:} Việc một cặp thuốc vắng mặt trong DDInter hoặc DrugBank không đồng nghĩa với âm tính tuyệt đối; do đó chỉ số đặc hiệu trên dữ liệu ngoài không được hỗ trợ.
  \item \textbf{An toàn Phần mềm và Kết cục Lâm sàng:} Nghiên cứu này chứng minh tính an toàn thông tin và ngăn chặn báo an toàn sai không có căn cứ ở cấp độ kiểm thử phần mềm; chưa phải là thử nghiệm ngẫu nhiên có đối chứng về giảm biến cố có hại do thuốc (ADE) trên bệnh nhân thực tế~\cite{holbrook2025alerts}.
\end{enumerate}

\section{Thảo luận và Ý nghĩa Lâm sàng}
Trọng tâm cốt lõi của CareGuard-VN nằm ở sự phân biệt rành mạch: \emph{kho tri thức DDI có thể hoàn toàn chính xác, nhưng kết luận y tế sẽ sai lệch nghiêm trọng nếu định danh hoạt chất đầu vào bị nhầm lẫn}. Việc tách rời chuẩn hóa khỏi kiểm tra tương tác tạo ra điểm mù nguy hiểm khi sai số chuẩn hóa bị che giấu dưới nhãn ``An toàn''.

Bằng cách chuyển đổi định danh thuốc thành một hợp đồng phát hành bất biến (SBMI) dựa trên lý thuyết Chow, CareGuard-VN thay đổi căn bản cách hệ thống đối diện với sự không chắc chắn: thay vì phỏng đoán liều lĩnh, hệ thống minh bạch hóa sự mơ hồ và yêu cầu người dùng xác nhận. Trên bộ thử nghiệm thực nghiệm nội bộ, CareGuard-VN đạt độ thu hồi 100,0\% trên các thử thách an toàn thuốc, độ tuân thủ 96,80\% trên đồ thị tri thức DrugBank 5.0 và tỷ lệ chặn fail-closed 100,0\% qua rào chắn FIDES, đồng thời đặt ra khung giao thức chuẩn mực cho đánh giá ngoài trong tương lai.

\section{Kết luận}
CareGuard-VN chứng minh rằng việc biến định danh thuốc gắn với nguồn thành điều kiện tiên quyết bắt buộc giúp ngăn chặn các báo cáo an toàn sai không có căn cứ trong các hệ thống CDSS kiểm tra tương tác thuốc. Dựa trên lý thuyết phân loại chọn lọc Chow và được kiểm chứng nội bộ trên các thử thách an toàn thuốc, kiểm chứng tuân thủ đồ thị tri thức (242/250, 96,80\% [93,81\%, 98,37\%]) và rào chắn an toàn fail-closed FIDES (100,0\% [95,86\%, 100,00\%]), CareGuard-VN thiết lập một chuẩn mực kiến trúc an toàn vững chắc. Công việc tiếp theo sẽ thực thi giao thức đánh giá ngoài (CG-01--CG-07) ngay khi tiếp nhận dữ liệu đăng ký thuốc quốc gia chính thức.

\section*{Tuyên bố Tính sẵn sàng của Dữ liệu và Mã nguồn}
Toàn bộ mã nguồn, kịch bản thử nghiệm và công cụ tái lập được công khai tại kho mã nguồn CLARA-Care: \repo. Dữ liệu DrugBank XML tuân thủ các điều khoản bản quyền tương ứng.

\section*{Đạo đức và Cam kết}
Nghiên cứu đánh giá hành vi an toàn phần mềm và các ca kiểm thử ẩn danh hóa. Các tác giả cam kết không có xung đột lợi ích.

\begin{thebibliography}{25}

\bibitem{nelson2011rxnorm}
S.~J. Nelson, K. Zeng, J. Kilbourne, T. Powell, and R. Moore,
``Normalized names for clinical drugs: RxNorm at 6 years,''
\emph{J. Am. Med. Inform. Assoc.}, vol. 18, no. 4, pp. 441--448, 2011.

\bibitem{sohn2014medxn}
S. Sohn, C. Clark, S. R. Halgrim, S. P. Murphy, C. G. Chute, and H. Liu,
``MedXN: an open source medication extraction and normalization tool for clinical text,''
\emph{J. Am. Med. Inform. Assoc.}, vol. 21, no. 5, pp. 858--865, 2014.

\bibitem{peters2011approx}
L. Peters, D. T. Nguyen, O. Bodenreider, and J. Kapusnik-Uner,
``An approximate matching method for clinical drug names,''
\emph{AMIA Annu. Symp. Proc.}, 2011.

\bibitem{newmangriffis2021ambiguity}
D. Newman-Griffis, G. Divita, B. Desmet, et al.,
``Ambiguity in medical concept normalization: An analysis of types and coverage in electronic health record datasets,''
\emph{J. Am. Med. Inform. Assoc.}, vol. 28, no. 3, pp. 516--532, 2021.

\bibitem{swaminathan2024selective}
A. Swaminathan, I. Lopez, W. Wang, et al.,
``Selective prediction for extracting unstructured clinical data,''
\emph{J. Am. Med. Inform. Assoc.}, vol. 31, no. 1, pp. 188--197, 2024.

\bibitem{chow1970}
C.~K. Chow,
``On optimum recognition error and reject tradeoff,''
\emph{IEEE Trans. Inf. Theory}, vol. 16, no. 1, pp. 41--46, 1970.

\bibitem{geifman2017selective}
Y. Geifman and R. El-Yaniv,
``Selective classification for deep neural networks,''
in \emph{Advances in Neural Information Processing Systems (NeurIPS)}, vol. 30, pp. 4878--4887, 2017.

\bibitem{gallifant2024rabbits}
J. Gallifant, S. Chen, P. J. F. Moreira, et al.,
``Language Models are Surprisingly Fragile to Drug Names in Biomedical Benchmarks,''
in \emph{Findings of EMNLP 2024}, pp. 12448--12465, 2024.

\bibitem{chen2024medguard}
C.-Y. Chen, Y.-L. Chen, J. Scholl, H.-C. Yang, and Y.-C. J. Li,
``Ability of machine-learning based clinical decision support system to reduce alert fatigue, wrong-drug errors, and alert users about look alike, sound alike medication,''
\emph{Comput. Methods Programs Biomed.}, vol. 243, 107869, 2024.

\bibitem{gibson2026ddireview}
B. N. Gibson, S. Chodapuneedi, H. Y. Koh, et al.,
``Identification, reliability, and validity of drug--drug interaction checkers in chronic diseases: A systematic review,''
\emph{Br. J. Pharmacol.}, vol. 183, no. 16, pp. 4532--4561, 2026.

\bibitem{holbrook2025alerts}
A. M. Holbrook, J. M. Silva, J. A. Y. Faruque, J. Deng, T. Schneider, and A. Jaffer,
``Effect of electronic drug--drug interaction alerts on patient and clinician outcomes: a systematic review,''
\emph{J. Am. Med. Inform. Assoc.}, vol. 32, no. 10, pp. 1617--1628, 2025.

\bibitem{wishart2018drugbank}
D. S. Wishart et al., ``DrugBank 5.0: a major update to the DrugBank database for 2018,''
\emph{Nucleic Acids Res.}, vol. 46, no. D1, pp. D1074--D1082, 2018.

\bibitem{tian2025ddinter}
Y. Tian et al., ``DDInter 2.0: an enhanced drug interaction resource with expanded data coverage, new interaction types, and improved user interface,''
\emph{Nucleic Acids Res.}, vol. 53, no. D1, pp. D1356--D1362, 2025.

\bibitem{ddinterdownload2026}
DDInter 2.0, ``Download portal,'' 2026. [Online]. Available: \url{https://ddinter2.scbdd.com/download/}.

\bibitem{korpela2026rxmap}
E. Korpela, L. H. Rubin, R. M. Dastgheyb, and Y. Xu,
``RxMap: an LLM-assisted tool for medication normalization,''
\emph{JAMIA Open}, vol. 9, no. 3, ooag085, 2026.

\bibitem{huser2025rxembed}
M. H\"user, J. Doole, V. Pinho, et al.,
``A machine learning approach for automating review of a RxNorm medication mapping pipeline output,''
\emph{J. Biomed. Inform.}, vol. 170, 104909, 2025.

\bibitem{crossddi2026}
H. Wang, B. Chu, and N. Fuhr, ``CrossDDI: Cross-Source Evidence-Grounded Drug-Drug Interaction Verification,'' in \emph{Proc. BioNLP 2026 Workshop}, Association for Computational Linguistics, 2026.

\bibitem{chase2026ddi}
A. Chase, A. Most, A. Sikora, et al.,
``Evaluation of large language models' ability to identify clinically relevant drug--drug interactions and generate high-quality clinical pharmacotherapy recommendations,''
\emph{Am. J. Health-Syst. Pharm.}, vol. 83, no. 13, pp. e494--e500, 2026.

\bibitem{nlmrxnorm2026}
U.S. National Library of Medicine, ``RxNorm Files: July 6, 2026 Monthly Release and Current Prescribable Content,'' 2026.

\bibitem{dailymed2026}
U.S. National Library of Medicine, ``DailyMed Structured Product Labeling Resources,'' 2026.

\bibitem{dav2026}
Drug Administration of Vietnam, Ministry of Health, ``Public drug registration and over-the-counter medicine lookup portal,'' 2026.

\bibitem{luo2020n2c2}
Y.-F. Luo, W. Sun, and A. Rumshisky, ``The 2019 n2c2/UMass Lowell shared task on clinical concept normalization,'' \emph{J. Am. Med. Inform. Assoc.}, 2020.

\bibitem{chen2026mcnsurvey}
H. Chen, Y. Zhou, R. Li, A. M. Illa, A. Cleveland, and J. Ding, ``A comprehensive survey on medical concept normalization: Datasets, techniques, applications, and future directions,'' \emph{J. Biomed. Inform.}, vol. 178, 105005, 2026.

\end{thebibliography}
\end{document}
